\documentclass{article}
\usepackage[margin=1.5cm]{geometry}
\usepackage{amsmath}
\usepackage{enumitem}

\setlist{nosep,leftmargin=*}

\begin{document}
\twocolumn

\title{Chapter 11 Exercise: From Classes to Modules}
\author{Prof. Jordan C. Hanson}
\date{}

\maketitle

\section{Build the Module}

Create a Python file named

\begin{verbatim}
complex_number.py
\end{verbatim}

Begin with the class from the warm-up:

\begin{verbatim}
class ComplexNumber:
    def __init__(self, real, imag):
        self.real = real
        self.imag = imag
    def __add__(self, other):
        return ComplexNumber(
            self.real + other.real,
            self.imag + other.imag)
    def __sub__(self, other):
        return ComplexNumber(
            self.real - other.real,
            self.imag - other.imag)
\end{verbatim}

Each object represents

\[
z=a+bi.
\]

\section{Complex Multiplication}

For

\[
z_1=a+bi,\qquad z_2=c+di,
\]

multiplication gives

\[
z_1z_2
=(ac-bd)+(ad+bc)i.
\]

Add this method:

\begin{verbatim}
def __mul__(self, other):
    real = self.real*other.real \
         - self.imag*other.imag
    imag = self.real*other.imag \
         + self.imag*other.real
    return ComplexNumber(real, imag)
\end{verbatim}

\section{Complex Conjugation}

The \textbf{complex conjugate} changes the sign of the
imaginary part:

\[
z=a+bi
\qquad\Longrightarrow\qquad
z^*=a-bi.
\]

Add

\begin{verbatim}
def conjugate(self):
    return ComplexNumber(
        self.real, -self.imag)
\end{verbatim}

\section{Complex Division}

To divide complex numbers, use the following formula:

\[
\frac{a+bi}{c+di}
=
\frac{ac+bd}{c^2+d^2}
+
\frac{bc-ad}{c^2+d^2}i.
\]

Add

\begin{verbatim}
def __truediv__(self, other):
    den = other.real**2 \
        + other.imag**2
    real = (self.real*other.real \
          + self.imag*other.imag)/den
    imag = (self.imag*other.real \
          - self.real*other.imag)/den
    return ComplexNumber(real, imag)
\end{verbatim}

\section{Create \texttt{main.py}}

In a second file, create

\begin{verbatim}
from complex_number import ComplexNumber
z1 = ComplexNumber(3, 2)
z2 = ComplexNumber(5, -4)
z_add = z1 + z2
z_sub = z1 - z2
z_mul = z1 * z2
z_div = z1 / z2
z_conj = z1.conjugate()
\end{verbatim}

Print the real and imaginary parts:

\begin{verbatim}
print("Add:",
      z_add.real, z_add.imag)
print("Subtract:",
      z_sub.real, z_sub.imag)
print("Multiply:",
      z_mul.real, z_mul.imag)
print("Divide:",
      z_div.real, z_div.imag)
print("Conjugate:",
      z_conj.real, z_conj.imag)
\end{verbatim}

Compare your hand calculations with the values printed by
\texttt{main.py}.

\end{document}